
\documentclass[10pt,openany,a4paper]{article}
\usepackage{tikz, pgfplots, tkz-euclide,calc}
\usetikzlibrary{patterns,shapes.arrows,calc}
\usepackage{fancyhdr}
\usepackage{enumerate,enumitem}
\usepackage{cancel}
\usepackage{varwidth}

% TAMBAHKAN PACKAGE SENDIRI KALAU KURANG

\usepackage{geometry}
\geometry{
	left = 22mm,
	right = 22mm,
	top = 25mm,
	bottom = 25mm,
}
\usepackage{graphicx}
\usepackage{subcaption}
\usepackage{soul}
\usepackage{hyperref}
\usepackage{float}
\usepackage{setspace}
\usepackage{afterpage}
\usepackage{xcolor}
\usepackage[nosectionbib]{apacite}
\usepackage{lmodern}
\usepackage{mathtools}
\usepackage{upgreek}
\usepackage{siunitx}
\usepackage{physics}
\usepackage{hyperref}
\renewcommand{\baselinestretch}{1.3}
\usepackage{multicol}
\usepackage{hyperref,ragged2e,amsmath,multicol,gensymb,setspace,
fancyhdr,amsfonts,tikz,pgfplots,nccmath,enumerate,verbatim,amsthm,physics}
\usepgfplotslibrary{polar,fillbetween}
% Disable tikz externalization to avoid requiring -shell-escape
\usepgfplotslibrary{external}
\tikzexternalize[prefix=tikz/]
\tikzset{reuse path/.code={\pgfsyssoftpath@setcurrentpath{#1}}}
\usetikzlibrary{shapes.geometric,arrows}
\newcommand\mylog[1]{\mathop{{}^{#1}\mathrm{log}}}
\pgfplotsset{compat=1.15}
\pgfplotsset{my style/.append style={axis x line=middle, axis y line=
middle, xlabel={$x$}, ylabel={$y$}, axis equal }}
\usepackage{etoolbox}
% remove stray \makeatother
\newtheorem{defn}{Definisi}[section]
\newtheorem{teo}[defn]{Teorema}
\newtheorem{thm}{Teorema}[section]
\newtheorem{lemma}{Lemma}
\newtheorem{lemmas}[defn]{Lemma}
\newtheorem{cor}[defn]{Akibat}
\newtheorem{asumsi}[defn]{Asumsi}
\theoremstyle{definition}
\newtheorem{con}{Contoh}[section]
\theoremstyle{definition}
\theoremstyle{plain}
\newtheorem{prop}{Proposisi}[section]
\renewcommand{\proofname}{Bukti}
\renewcommand{\thethm}{\arabic{chapter}.\arabic{thm}}

\newcommand{\normlp}[1]{\left\|#1\right\|_{L^p_{\lambda}(0,T;H)}}% Fungsi norm (||x||)
\newcommand{\R}{\mathbb{R}}
\renewcommand{\natural}{\mathbb{N}}
\newcommand{\Xpl}{X^R_{p,\lambda}}
\newcommand{\Xpls}{X^R_{p,\lambda^*}}
\newcommand{\qpl}{q^R_{p,\lambda}}
\newcommand{\qpls}{q^R_{p,\lambda^*}}
\pagestyle{fancy}
\fancyhf{}
\lhead{Halaman \thepage}
\rhead{Pembahasan Soal EAS 2024/2025 \\ (\href{https://instagram.com/ahmadzakiyudin_/}{@ahmadzakiyudin\_})}
\hypersetup{
    colorlinks=true,
    linkcolor=blue,
    filecolor=blue,      
    urlcolor=blue,
}
\newenvironment{jawab}[2][tes]{
% \fbox{\begin{minipage}{14.65cm}
\textbf{Kata kunci materi:} #1.\\
\textbf{Penyelesaian:} #2
% \end{minipage}}
}{}

\newenvironment{jawab2}[2][tes]{
\fbox{\begin{minipage}{14.65cm}
 #2
\end{minipage}}}{}

\newenvironment{jawaban}[2][tes]{
\colorlet{oldcolor}{.}
        \color{red}
        \setlength{\fboxrule}{1pt}\boxed{\color{oldcolor}#2}\color{oldcolor}.
}{}

\newenvironment{jawaban2}[2][tes]{
\colorlet{oldcolor}{.}
        \color{red}
        \setlength{\fboxrule}{1pt}\fbox{\color{oldcolor}#2}\color{oldcolor}.
}{}

\newenvironment{penting}[2][tes]{
\colorlet{oldcolor}{.}
\color{blue}
\fbox{\color{oldcolor}#2}\color{oldcolor}
}{}

\begin{document}
% \begin{titlepage}
%     \vspace*{\fill}
%     \begin{center}
%       \Huge {PEMBAHASAN SOAL EAS \\ KALKULUS II \\ TAHUN 2024/2025}\\[0.4 cm]
%       \huge {Ahmad Hisbu Zakiyudin}\\[0.2 cm]
%       \Large {(\href{https://instagram.com/ahmadzakiyudin_/}{@ahmadzakiyudin\_})}
%     \end{center}
%     \vspace*{\fill}
%   \end{titlepage}
% !TEX root = ../../BookTemplate.tex
%%%%%%%%%%%%%%%%%%%%%%%%%%%%%%%%%%%%%%%%%%%%%%%%%%%%%%%%%%%%%%%%%%%%%%%%%%%%%%%%%%
% \begin{titlepage}
%     \raggedleft	
% 	\hspace{.025\textwidth}
% 	\parbox[b]{\textwidth}{
%     \vspace{1.5cm}
% 		{\Huge\bfseries Pembahasan Soal EAS}                \\[20pt]
% 		{\huge\textit\ Kalkulus II}\\ \\[10pt]
%         {\Large\textsc{Tahun 2024/2025}}}
% \vfill
% \rule{1pt}{.14\textheight}
% \hspace{.025\textwidth}
%     \raisebox{10pt}{\parbox[b]{.95\textwidth}{ 
%         {\Large\textsc{Ahmad Hisbu Zakiyudin}}              \\[10pt]
%         \large{\href{https://instagram.com/ahmadzakiyudin_/}{@ahmadzakiyudin\_}}               \\[25pt]
%         {\large 5 Juni 2026}}}
% \end{titlepage}
\begin{titlepage}
	\centering % Center everything on the title page
	% \scshape % Use small caps for all text on the title page
	\vspace*{1.5\baselineskip} % White space at the top of the page
% ===================
%	Title Section 	
% ===================

	\rule{13cm}{1.6pt}\vspace*{-\baselineskip}\vspace*{2pt} % Thick horizontal rule
	\rule{13cm}{0.4pt} % Thin horizontal rule
	
		\vspace{0.75\baselineskip} % Whitespace above the title
% ========== Title ===============	
	{	\Huge Pembahasan Soal EAS\\ 
			\vspace{4mm}
		Kalkulus II \\ \vspace{2mm} \huge{Tahun 2024/2025}\\	}
% ======================================
		\vspace{0.75\baselineskip} % Whitespace below the title
	\rule{13cm}{0.4pt}\vspace*{-\baselineskip}\vspace{3.2pt} % Thin horizontal rule
	\rule{13cm}{1.6pt} % Thick horizontal rule
	
		\vspace{1.75\baselineskip} % Whitespace after the title block
% =================
%	Information	
% =================
{\large Oleh: Ahmad Hisbu Zakiyudin}
	% {\large Oleh: Ahmad Hisbu Zakiyudin \\
	% 	\vspace*{1.2\baselineskip}
	% \href{https://instagram.com/ahmadzakiyudin_/}{@ahmadzakiyudin\_}} \\
	\vfill
	\vspace{1mm}
\huge{Selamat Belajar!}
\end{titlepage}
\makeatletter
\setlength\fboxsep{3pt}
% ensure enough header height for fancyhdr
\setlength{\headheight}{24pt}
\section*{Kelas 29, 30, 37, 67, 68}
\begin{enumerate}
    \item Sketch the region bounded by $y=\sqrt{x}$, $y=6-x$, and the $x$-axis. Then find the area of the region.
    \item Find the centroid of the homogeneous region bounded by $y=2x$, $y=x+1$, and the $y$-axis.
    \item Find $\dfrac{dx}{dt}$ and $\dfrac{dy}{dt}$ for the parametric equations $x=2t^3-15t^2+24t+7$ and $y=t^2+t+1$, Then determine all values of $t$ for which the tangent line is vertical.
    \item Sketch the curve $r=2-2\cos\theta$, and find the area inside the curve.
    \item Consider the sequence
    \begin{align*}
        \left\{\frac{n^2(1/2)^n}{2n^2}\right\}_{n=1}^{\infty}.
    \end{align*}
    Write the first three terms and determine whether the sequence is converges or diverges.
\end{enumerate}
\section*{Kelas 61-65}
\begin{enumerate}
    \item Gambarkan daerah yang dibatasi oleh $y=x^2+1, y=3-x, x=0$, dan sumbu-$x$, dan kemudian dapatkan luas daerah tersebut.
    \item Gunakan Dalil Guldin untuk mendapatkan volume benda yang diperoleh dengan memutar lingkaran $x^2-6x+y^2-4y+12=0$ terhadap sumbu-$x$.
    \item Diketahui persamaan parametrik $x=3+2\cos t$ dan $y=2+2\sin t$. Gunakan identitas trigonometri yang sesuai untuk mengeliminasi parameter $t$ dan kemudian gambarkan grafiknya.
    \item Gambarkan kurva dari rose $r=8\sin2\theta$, dan kemudian dapatkan luas daerah di dalam kurva tersebut.
    \item Diketahui barisan
    \begin{align*}
        \left\{\frac{n^2-6n+20}{n+1}\right\}_{n=1}^{\infty}.
    \end{align*}
    Dapatkan $a_n$ dan $a_{n+1}$. Kemudian, dengan uji selisih, carilah bilangan bulat $N$ sedemikian sehingga barisan tersebut adalah monoton untuk $n\geq N$.
\end{enumerate}
\section*{Kelas 15-20 dan 22-28}
\begin{enumerate}
    \item Gambarkan daerah yang dibatasi oleh $y=\sqrt{x+4}, y=2-x, x=1$, dan sumbu-$x$, dan kemudian dapatkan luas daerah tersebut.
    \item Dapatkan titik berat dari bidang datar homogen yang dibatasi oleh $y=x^2-2x-1$ dan $y=2$.
    \item Dapatkan $\dfrac{dy}{dx}$ dari persamaan parametrik $x=\cos 2t$ dan $y=\sin 2t$, $0\leq t \leq \dfrac{\pi}{2}$. Selanjutnya dapatkan panjang busurnya
    \item Gambarkan kurva dari rose $r=4\sin3\theta$, dan kemudian dapatkan luas daerah di dalam kurva tersebut.
    \item Diketahui barisan 
    \begin{align*}
        \left\{\frac{n^2-8n+30}{n+1}\right\}_{n=1}^{\infty}.
    \end{align*}
    Dapatkan $a_n$ dan $a_{n+1}$. Kemudian, dengan uji selisih, carilah bilangan bulat $N$ sedemikian sehingga barisan tersebut adalah monoton untuk $n\geq N$.
\end{enumerate}
\section*{Kelas 51-60}
\begin{enumerate}
    \item Gambarkan daerah yang dibatasi oleh $y=x^2, y=x+2, x=0$, dan $x=1$. Dapatkan volume benda padat jika daerah tersebut diputar terhadap garis $y=-1$.\\
    Catatan: perhatikan jarak antara kurva dengan sumbu putar.
    \item Dapatkan turunan dari $f(x)=\sqrt{1-x^2}$ dan kemudian gunakan integral untuk mendapatkan panjang busurnya dari $x=0$ sampai $x=1$.
    Catatan: $\displaystyle \int \frac{1}{\sqrt{1-x^2}}dx=\sin^{-1} x $.
    \item Dapatkan semua nilai $\theta$ yang menyebabkan $r=4+4\cos \theta$ mempunyai garis singgung horizontal.
    \item Gambarkan daerah di dalam lingkaran $r=\sqrt{3}\sin\theta$ dan di luar kardioida $r=1+\cos\theta$, dan kemudian dapatkan luas daerah tersebut.
    \item Tentukan nilai dari $\displaystyle \sum_{k=1}^{\infty} \frac{2}{k^2+2k}$.
    \end{enumerate}
\vfill
\vspace{1mm}
\begin{center}
     Jika anda menemukan kesalahan pada pembahasan, 
atau membutuhkan tutor untuk mata kuliah Kalkulus, hubungi saya melalui Instagram \href{https://instagram.com/ahmadzakiyudin_/}{@ahmadzakiyudin\_} 
\end{center}
\newpage

\end{document}


